\documentclass[10pt]{article}
\usepackage[margin=0.72in]{geometry}
\usepackage{longtable}
\usepackage{array}
\usepackage{xcolor}
\usepackage{hyperref}
\definecolor{accent}{HTML}{173B57}
\definecolor{redline}{HTML}{8B1E1E}
\newcommand{\code}[1]{\texttt{\small #1}}
\newcommand{\pathline}[1]{\par\noindent\texttt{\scriptsize\detokenize{#1}}\par}
\title{AJCR Horizon Status and Route-Divergence Brief}
\author{Human mathematician handoff}
\date{24 August 2026}
\begin{document}
\maketitle

\noindent\textbf{Snapshot:} 2026-08-24 10:06 +0800.\quad
\textbf{Source ledger:} \code{a58a86aca1}.

\section*{Purpose and verdict}
This is a status handoff against the two supplied review PDFs. It records what Horizon has completed relative to the reviewed route, what remains, and why repeated sessions have not crossed the same Phase 7 boundary.

The rank-one route has reached the separably-closed milestone:
\begin{verbatim}
PicRankOneOpen / DivRankOneOpen
  -> canonicalRankOneAbelIso
  -> translated rank-one cover over a separably closed field
  -> pic0_sepClosed_representableBy
\end{verbatim}
The current blocker is finite-Galois descent, not the old family-level rank-one producer. Horizon is pursuing:
\begin{verbatim}
finite-stage atlas and GlueData
  -> P.gluedOver and glued base-change maps
  -> universal class and overlap/glue equivalences
  -> binder-free RepresentableBy P.gluedOver
  -> exact-carrier FiniteInAffine / projectivity
  -> original-field pic0_representableBy
  -> PicRepDatum -> JacobianData -> Challenge.Jacobian
\end{verbatim}
The first line is infrastructure. The remaining lines contain three missing mathematical producers and a resource-bound compile frontier.

\section*{Reviewed route versus current status}
The binding route in the review plan is:
\begin{verbatim}
PicRankOneOpen -> rankOneAbelIso
  -> rankOne_translate_cover_sepClosed
  -> pic0_sepClosed_representableBy
  -> finiteGaloisDescent -> pic0_representableBy -> JacobianData
\end{verbatim}

\begin{longtable}{>{\raggedright\arraybackslash}p{0.24\linewidth} >{\raggedright\arraybackslash}p{0.18\linewidth} p{0.49\linewidth}}
\textbf{Phase} & \textbf{Board} & \textbf{Measured meaning}\\\hline
P3 rank-one loci & Done & Public \code{PicRankOneOpen}/\code{DivRankOneOpen}, openness, and base change are reported unconditional.\\
P4 rank-one Abel iso & Done & \code{canonicalRankOneAbelIso} is root-reachable and consumed by the separably-closed route; this is the family-level producer emphasized by the reviewer.\\
P5 translated cover & Done & The translator is tied to the input Picard class and lands in the public rank-one locus.\\
P6 sep-closed Pic0 & Done & \code{pic0\_sepClosed\_representableBy} and same-carrier \code{picRepDatumSepClosed}/\code{jacobianDataSepClosed} are present.\\
P7 finite-Galois descent & Blocked & No unconditional original-field producer. The finite-stage glue cone is only partly native-built; universal Yoneda and exact-carrier orbit gates remain open.\\
P8 Jacobian & Blocked & No original-field \code{J}/\code{rep} pair reaches \code{PicRepDatum}, \code{JacobianData}, and \code{Challenge.Jacobian}.\\
\end{longtable}

\section*{Remaining gates}
\begin{longtable}{>{\raggedright\arraybackslash}p{0.25\linewidth} >{\raggedright\arraybackslash}p{0.16\linewidth} p{0.48\linewidth}}
\textbf{Gate} & \textbf{Kind} & \textbf{Closure condition}\\\hline
Native glue root & Build & Fresh root import certifies \code{GluingDiagramIso}, top \code{PreSnd}, \code{OverlapIsoSnd}, and \code{GluedComparison}. \code{GluePackage.olean} appeared on 24 Aug 00:29; these four downstream artifacts remain absent.\\
Finite-stage representation & Math & A theorem of the form \code{pic0RepresentableBy\_finiteStageGlue} returns \code{RepresentableBy P.gluedOver}, with a natural \code{homEquiv} and no \code{(rep : ...)} binder.\\
Exact-carrier orbit geometry & Math & Prove \code{P.gluedMap.IsProjective}, \code{FiniteInAffine P.glueData.glued}, or an equivalent result for the same carrier. Existing orbit lemmas consume this hypothesis.\\
Original-field Pic0 & Math & A descent theorem returns both the scheme and its \code{RepresentableBy} certificate, immediately consumed by \code{pic0\_representableBy}.\\
Jacobian handoff & Integration & Feed the same original-field \code{J}/\code{rep} through \code{PicRepDatum} and \code{JacobianData}, then into \code{Challenge.lean} without an import cycle.\\
\end{longtable}

\section*{Why the same issues recur}
\begin{enumerate}
\item \textbf{The post-P6 route is over-decomposed.} Horizon split one producer-level Phase 7 edge into transition models, triple overlaps, ring maps, glue records, base-change isomorphisms, overlap projections, and only later the Yoneda equivalence. Prerequisites create commits but do not close an acceptance edge.
\item \textbf{Source labels are over-read.} The graph marks some declarations \code{lean\_ok} even when their native artifacts are absent. \code{Pic0CriticalPath.lean} imports the top glue modules, so only a native artifact plus fresh root import and axiom print is evidence.
\item \textbf{The compile cone consumes the iteration budget.} The GluePackage probe used about 7.1 GB RSS. The first expensive declaration in the next cone is \code{overlapBaseChangeIso\_hom\_iota} at \code{Pic0FiniteStageGluingDiagramIso.lean:276}; historical checks reached about 610 seconds.
\item \textbf{Orchestration leaves dead work looking active.} Run 0154 metadata still says \code{running} from \code{2026-08-23T14:55:00Z}, while \code{horizon ps} reports a zombie marker and no live Lean process. Its system sessions returned the task to \code{queued}. Runs 0145 and 0148 repeatedly relaunched the same alignment task; later sessions mostly returned \code{exit 1} or \code{queued}.
\item \textbf{The custom carrier is an architectural question.} The review plan asks to reuse the AJC finite-Galois quotient/gluing engine. The current route first builds \code{P.gluedOver} and then asks for a new universal \code{homEquiv}. This may be right, but it may also be a second descent stack without a proved identification with the required Pic0 functor.
\end{enumerate}

\section*{Run chronology}
\begin{longtable}{>{\raggedright\arraybackslash}p{0.12\linewidth} p{0.16\linewidth} p{0.54\linewidth}}
\textbf{Run} & \textbf{Date} & \textbf{Result}\\\hline
0145 & 14 Aug & Eight-round alignment task; later sessions repeatedly returned \code{exit 1}/\code{queued}.\\
0148 & 14 Aug & Same alignment task relaunched; metadata/audit work, then queued or failed sessions.\\
0149 & 14 Aug onward & Audit classified 266 commits as 0 acceptance edges, 82 consumed prerequisites, and 10 conditional consumers.\\
0153 & 22 Aug & Compile isolation measured the GluePackage/DiagramIso frontier.\\
0154 & 23--24 Aug & Compile repair produced \code{GluePackage.olean}, but no root-certified downstream top cone; stale \code{running} metadata remains.\\
0156 & 23 Aug & Roadmap decomposition separated build, Yoneda, orbit, and original-field gates; no theorem edge was added.\\
\end{longtable}

\section*{Questions for review}
\begin{enumerate}
\item Does \code{P.gluedOver} implement the finite-level representer required by the review plan, including its universal class and natural Yoneda equivalence, or is it only an atlas/base-change model?
\item Can the existing AJC \code{representableBy\_of\_finiteGalois\_baseChange} theorem consume a smaller finite-level producer instead of extending this custom glue cone?
\item What mathematical input proves orbit affineness or projectivity for the exact glued carrier?
\item Is the missing binder-free \code{RepresentableBy} theorem a functor/sheaf descent argument rather than another overlap identity?
\item Which single theorem should be the next bounded acceptance milestone?
\end{enumerate}

\section*{Reading and evidence}
\pathline{MainProjects/Algebraic-Jacobian-Challenge/informal/Lean_Algebraic_Jacobian_Complete_Execution_Plan.pdf}
\pathline{MainProjects/Algebraic-Jacobian-Challenge/informal/AJCR_Runs_121_122_123_Supervision_Note_2026-08-07.pdf}
\pathline{MainProjects/Algebraic-Jacobian-Challenge-Rebuild/AlgebraicJacobian/Picard/Pic0CriticalPath.lean}
\pathline{.archon-horizon/runs/0152/sessions/0006-horizon-ajcr-strategy-review/report.md}
\pathline{.archon-horizon/runs/0153/sessions/0002-horizon-ajcr-compile-isolation/report.md}
\pathline{.archon-horizon/runs/0154/sessions/0002-horizon-ajcr-compile-frontier-repair/report.md}
\pathline{.archon-horizon/runs/0154/sessions/0005-horizon-ajcr-compile-frontier-repair/report.md}
\pathline{.archon-horizon/runs/0156/sessions/0002-horizon-ajcr-roadmap-decomposition/report.md}

The existing rebuild README is useful but its artifact inventory predates the latest \code{GluePackage.olean}; do not use that inventory as current. The task database currently has 28 done, 3 blocked, 4 queued, 1 running, 1 failed, and 45 cancelled tasks. The dashboard also has a blank \code{generatedAt} and stale Aug 21 status PDFs.

\bigskip
\noindent\textbf{Bottom line.} Horizon is not stuck because the rank-one route failed. It is stuck because Phase 7 interleaves a resource-bound custom glue cone with missing mathematical producers. A human decision on whether \code{P.gluedOver} is the right finite-level object, followed by one producer-level acceptance test, is needed before more sessions can be expected to move the headline theorem.
\end{document}
